\chapter{Multiplexers and Demultiplexers}
\label{ch:mux}

In Chapter~\ref{ch:coders} we learned to light \emph{one} line out of $2^n$ by
its number: a decoder answers the question ``which one is it?''. But a line by
itself carries nothing. A question remains: how can we use a number not just to
point at a signal but to \emph{pass it through}, that is, to choose one of
several sources and send it onward? The answer is the multiplexer, and we will
see that inside it is built from the decoder we already know.

\begin{buildbox}
\textbf{Task.} Choose whose value to show on an LED: switch \code{a} or switch
\code{b}, according to a third switch \code{s}.

\textbf{In Logra.} Learn tab, the ``Circuits and the board'' book, the
``Selection and distribution'' chapter, the lesson ``2-to-1 multiplexer''
(Designer). You need one ``2:1 multiplexer'' block from the Select group.

\textbf{Expected result.} When $s=0$ the LED repeats \code{a}; when $s=1$ it
repeats \code{b}. The rest of this chapter explains what is inside this block
and what it can be turned into.
\end{buildbox}

\section{A Digital Switch}

Picture an old rotary switch: you turn the knob and it connects a common
contact to one of several others. A \term{multiplexer} (MUX) does the same with
digital signals: it has several data inputs, one output, and \term{select}
(control) inputs that choose which data input is connected to the output.

\begin{figure}[H]
\centering
\begin{tikzpicture}
  % mechanical switch
  \foreach \i in {0,...,3} {
    \fill (0,1.5-\i) circle (2pt);
    \draw[wire] (-1,1.5-\i) node[left] {$x_\i$} -- (0,1.5-\i);
  }
  \fill (2,0) circle (2.5pt);
  \draw[very thick, FpgaOrange] (2,0) -- (0.08,0.5);
  \draw[wire] (2,0) -- (3,0) node[right] {$y$};
  \node[lbl] at (1.2,-2.3) {mechanical switch};
  % multiplexer
  \begin{scope}[xshift=7.5cm]
    \draw[thick, fill=orange!15] (0,2) -- (1.6,1.2) -- (1.6,-1.2) -- (0,-2) -- cycle;
    \node[font=\sffamily\bfseries] at (0.8,0) {MUX};
    \foreach \i in {0,...,3} {
      \draw[wire] (-1,1.5-\i) node[left] {$x_\i$} -- (0,1.5-\i);
      \node[font=\sffamily\footnotesize, anchor=west] at (0.05,1.5-\i) {\i};
    }
    \draw[wire] (1.6,0) -- (2.6,0) node[right] {$y$};
    \draw[wire, FpgaBlue] (0.6,-1.7) -- (0.6,-2.6) node[below] {$s_1$};
    \draw[wire, FpgaBlue] (1.1,-1.45) -- (1.1,-2.6) node[below] {$s_0$};
    \node[lbl] at (0.8,-3.3) {digital 4:1 multiplexer};
  \end{scope}
  \draw[-{Stealth[length=4mm]}, line width=2pt, FpgaOrange] (4.2,0) -- (5.6,0);
\end{tikzpicture}
\caption{A multiplexer is a controlled switch: when $s_1s_0 = 01$ the output is
connected to input $x_1$}
\end{figure}

A multiplexer with $n$ select inputs can choose from $2^n$ data inputs. It is
written as a ratio: \textbf{2:1}, \textbf{4:1}, \textbf{8:1} (how many inputs :
how many outputs).

\begin{diybox}[A four-input switch]
\textbf{Task.} \stars{1} Repeat: in the Designer place a ``4:1 multiplexer''
block (the ``Number of inputs'' parameter is 4), connect four different sources
(switches or constants) to the inputs $D_0\ldots D_3$, and make the $S$ input
two bits wide. \stars{2} Modify: make input number 2 go to the output, and write
down which code is needed on $S$.

\textbf{In Logra.} The Designer and Waveforms tabs; the reference is the example
\code{mux-four}.

\textbf{Expected result.} For $S=00,01,10,11$ the output equals $D_0$, $D_1$,
$D_2$, $D_3$; the input number equals the binary number on $S$.

\textbf{How to check.} On the Waveforms tab go through the four values of $S$
and compare with Fig.~\ref{fig:mux41} and the timing diagram later in the text.
\end{diybox}

\section{The 2:1 Multiplexer}

The simplest multiplexer has two inputs $a$, $b$ and one select input $s$: if
$s = 0$, then $a$ passes to the output; if $s = 1$, then $b$ passes.

\begin{figure}[H]
\centering
\begin{minipage}{0.24\textwidth}
\centering
\begin{tikzpicture}
  \draw[thick, fill=orange!15] (0,1.1) -- (1.1,0.6) -- (1.1,-0.6) -- (0,-1.1) -- cycle;
  \node[font=\sffamily\small] at (0.55,0) {MUX};
  \draw[wire] (-0.7,0.6) node[left] {$a$} -- (0,0.6); \node[font=\sffamily\footnotesize, anchor=west] at (0.02,0.6) {0};
  \draw[wire] (-0.7,-0.6) node[left] {$b$} -- (0,-0.6); \node[font=\sffamily\footnotesize, anchor=west] at (0.02,-0.6) {1};
  \draw[wire] (1.1,0) -- (1.7,0) node[right] {$y$};
  \draw[wire, FpgaBlue] (0.55,-0.85) -- (0.55,-1.5) node[below] {$s$};
\end{tikzpicture}
\end{minipage}
\begin{minipage}{0.21\textwidth}
\centering
\begin{tabular}{c|c}
\toprule
$s$ & $y$ \\
\midrule
0 & $a$ \\
1 & $b$ \\
\bottomrule
\end{tabular}

\medskip
$y = \overline{s}\,a + s\,b$
\end{minipage}
\begin{minipage}{0.52\textwidth}
\centering
\begin{tikzpicture}
  \node[gAND] (g1) at (3.7,1.2) {};
  \node[gAND] (g2) at (3.7,-1.2) {};
  \node[gOR] (o) at (6.0,0) {};
  \node[gNOT] (n) at (1.5,0) {};
  \draw[wire] (g1.input 1) -- (g1.input 1 -| 0,0) node[left] {$a$};
  \draw[wire] (g2.input 2) -- (g2.input 2 -| 0,0) node[left] {$b$};
  \draw[wire, FpgaBlue] (0,0) node[left, text=FpgaDark] {$s$} -- (0.5,0) coordinate (sj);
  \fill[FpgaBlue] (sj) circle (1.3pt);
  \draw[wire, FpgaBlue] (sj) -- (n.input);
  \draw[wire, FpgaBlue] (sj) |- (g2.input 1);
  \draw[wire] (n.output) -| ([xshift=-4mm]g1.input 2) -- (g1.input 2);
  \draw[wire] (g1.output) -| ([xshift=-4mm]o.input 1) -- (o.input 1);
  \draw[wire] (g2.output) -| ([xshift=-4mm]o.input 2) -- (o.input 2);
  \draw[wire] (o.output) -- ++(0.5,0) node[right] {$y$};
\end{tikzpicture}
\end{minipage}
\caption{The 2:1 multiplexer: symbol, table, formula, and gate-level circuit}
\label{fig:mux21}
\end{figure}

How the circuit works: when $s = 0$ the upper AND gate ``lets $a$ through'',
while the lower one is ``closed'' (it has a 0 on its input). When $s = 1$ it is
the other way around. The OR gate collects the result.

\begin{diybox}[A multiplexer from gates]
\textbf{Task.} \stars{1} Build a 2:1 multiplexer from NOT, AND, and OR gates as
in Fig.~\ref{fig:mux21}. \stars{2} Replace all three kinds of gates with NAND
gates (Chapter~\ref{ch:gates}) and compare the number of gates.

\textbf{In Logra.} Learn tab, the ``Selection and distribution'' chapter, the
lesson ``A multiplexer from NOT, AND, OR''; for comparison, the ready-made block
and the example \code{mux}.

\textbf{Expected result.} Your circuit and the ready-made block give the same
output for all eight combinations of $a,b,s$.

\textbf{How to check.} Waveforms: feed the same inputs to both circuits and
compare the outputs; RTL: compare your drawing with the drawing of the block.
\end{diybox}

\section{The 4:1 Multiplexer}

In a 4:1 multiplexer, two select inputs $s_1 s_0$ give the number of the input,
from 0 to 3.
\[
  y = \overline{s_1}\,\overline{s_0}\,x_0 + \overline{s_1}\,s_0\,x_1 + s_1\,\overline{s_0}\,x_2 + s_1\,s_0\,x_3.
\]
Look closely at the formula: the factors multiplying $x_i$ are exactly the
outputs of the 2$\to$4 \emph{decoder} from Chapter~\ref{ch:coders}! So a
multiplexer consists of an address decoder, four AND gates, and one OR gate.

\begin{figure}[H]
\centering
\begin{tikzpicture}
  \draw[thick, fill=FpgaLight] (0,-5.5) rectangle (1.6,0.8);
  \node[font=\sffamily\bfseries] at (0.8,-2.35) {DC};
  \draw[wire, FpgaBlue] (-0.8,-1.5) node[left, text=FpgaDark] {$s_1$} -- (0,-1.5);
  \draw[wire, FpgaBlue] (-0.8,-3.2) node[left, text=FpgaDark] {$s_0$} -- (0,-3.2);
  \foreach \i in {0,...,3} {
    \node[gAND] (g\i) at (5.0,-\i*1.5-0.3) {};
    \draw[wire] (1.6,0 |- g\i.input 2) -- (g\i.input 2);
    \node[font=\sffamily\footnotesize, anchor=south east, inner sep=1pt] at (1.6,0 |- g\i.input 2) {\i};
    \draw[wire] (g\i.input 1) -- ++(-1.8,0) node[left, font=\small] {$x_\i$};
  }
  \node[gOR, logic gate inputs=nnnn, minimum height=18mm] (o) at (8.9,-2.55) {}; % fig-ok: 4-input gate, 18 mm
  \draw[wire] (g0.output) -| ([xshift=-5mm]o.input 1) -- (o.input 1);
  \draw[wire] (g1.output) -| ([xshift=-10mm]o.input 2) -- (o.input 2);
  \draw[wire] (g2.output) -| ([xshift=-10mm]o.input 3) -- (o.input 3);
  \draw[wire] (g3.output) -| ([xshift=-5mm]o.input 4) -- (o.input 4);
  \draw[wire] (o.output) -- ++(0.7,0) node[right] {$y$};
\end{tikzpicture}
\caption{The 4:1 multiplexer = address decoder + AND gates + OR gate}
\label{fig:mux41}
\end{figure}

\begin{figure}[H]
\centering
\begin{tikztimingtable}[timing/slope=0.05, xscale=0.9, yscale=1.4, timing/font=\sffamily\small, timing/rowdist=3.2ex]
  $x_0$     & 16{1L 1H} \\
  $x_1$     & 8{2L 2H} \\
  $x_2$     & 32H \\
  $x_3$     & 32L \\
  $s_1s_0$  & 8D{00} 8D{01} 8D{10} 8D{11} \\
  $y$       & 4{1L 1H} 2{2L 2H} 8H 8L \\
\end{tikztimingtable}
\caption{The 4:1 multiplexer passes different input signals to the output in
turn}
\end{figure}

\begin{diybox}[The 4:1 multiplexer]
\textbf{Task.} \stars{1} Go through the lesson and build the 4:1 multiplexer
from a block. \stars{2} Build it again ``as in Fig.~\ref{fig:mux41}'': an
address decoder, four AND gates, one OR gate.

\textbf{In Logra.} Learn tab, the ``Selection and distribution'' chapter, the
lesson ``4-to-1 multiplexer''; for the second task, the lesson ``Decoder'' as a
source of the decoder.

\textbf{Expected result.} Both circuits give the same table: the output equals
the input whose number is set on $s_1s_0$.

\textbf{How to check.} Waveforms: go through all $4$ addresses for different
values of $x_0\ldots x_3$ (two sets are enough); RTL: the block and your
assembly should differ only in appearance, not in logic.
\end{diybox}

\section{A Tree of Multiplexers}

A large multiplexer can be built from small ones. A 4:1 multiplexer is three
2:1 multiplexers: the first level chooses by the low address bit, the second by
the high bit. An 8:1 multiplexer is a tree of seven 2:1 multiplexers in three
levels, and so on.

\newcommand{\muxsym}[3]{% name, x, y
  \draw[thick, fill=orange!15] (#2,#3+0.55) -- (#2+0.7,#3+0.3) -- (#2+0.7,#3-0.3) -- (#2,#3-0.55) -- cycle;
  \coordinate (#1-i0) at (#2,#3+0.3); \coordinate (#1-i1) at (#2,#3-0.3);
  \coordinate (#1-o) at (#2+0.7,#3); \coordinate (#1-s) at (#2+0.35,#3-0.42);}

\begin{figure}[H]
\centering
\begin{tikzpicture}
  \muxsym{m0}{0}{1}
  \muxsym{m1}{0}{-1}
  \muxsym{m2}{2}{0}
  \foreach \m/\a/\b in {m0/x_0/x_1, m1/x_2/x_3} {
    \draw[wire] (\m-i0) -- ++(-0.6,0) node[left] {$\a$};
    \draw[wire] (\m-i1) -- ++(-0.6,0) node[left] {$\b$};
  }
  \draw[wire] (m0-o) -- ++(0.4,0) |- (m2-i0);
  \draw[wire] (m1-o) -- ++(0.4,0) |- (m2-i1);
  \draw[wire] (m2-o) -- ++(0.6,0) node[right] {$y$};
  \draw[wire, FpgaBlue] (m0-s) -- (0.35,-0.575);
  \fill[FpgaBlue] (0.35,0) circle (1.3pt);
  \draw[wire, FpgaBlue] (0.35,0) -- (-0.9,0) node[left, text=FpgaDark] {$s_0$};
  \draw[wire, FpgaBlue] (m2-s) -- (2.35,-1.2) node[below, text=FpgaDark] {$s_1$};
  \node[lbl] at (0.35,-2.4) {level 1};
  \node[lbl] at (2.35,-2.4) {level 2};
\end{tikzpicture}
\caption{A 4:1 multiplexer made of three 2:1 multiplexers}
\label{fig:muxtree}
\end{figure}

\begin{diybox}[A tree of multiplexers]
\textbf{Task.} \stars{1} Build a 4:1 multiplexer from three ``2:1 multiplexer''
blocks as in Fig.~\ref{fig:muxtree}. \stars{2} Add four more blocks to get an
8:1 multiplexer; write down which address bits go to which level.

\textbf{In Logra.} The Designer and RTL tabs; the reference is the example
\code{mux-four} (compare with the ready-made block).

\textbf{Expected result.} The tree and the ready-made block give the same
output for all addresses; the low address bit $s_0$ goes to the first level.

\textbf{How to check.} Waveforms: apply the constant values $0,1,0,1$ to
$x_0\ldots x_3$ and go through the addresses---the output should alternate
$0,1,0,1$.
\end{diybox}

\section{The Multiplexer as a Universal Gate}

Here is a remarkable property of the multiplexer. Apply the \emph{variables}
$a$ and $b$ to the select inputs, and apply the \emph{constants} 0 and 1,
taken from the right-hand column of some function's truth table, to the data
inputs. The multiplexer will then pick the needed row of the table and output
the value of the function! Any function of $n$ variables can be implemented with
a single $2^n$:1 multiplexer.

\begin{figure}[H]
\centering
\begin{tikzpicture}
  \draw[thick, fill=orange!15] (0,2) -- (1.6,1.2) -- (1.6,-1.2) -- (0,-2) -- cycle;
  \node[font=\sffamily\bfseries] at (0.8,0) {MUX};
  \foreach \i/\v/\ab in {0/0/00, 1/1/01, 2/1/10, 3/0/11} {
    \draw[wire] (-0.8,1.5-\i) -- (0,1.5-\i);
    \node[draw=FpgaGreen, fill=green!8, thick, minimum size=6mm, font=\ttfamily] at (-1.1,1.5-\i) {\v};
    \node[font=\sffamily\footnotesize, anchor=west] at (0.05,1.5-\i) {\i};
    \node[lbl, anchor=east] at (-1.5,1.5-\i) {$ab=\ab$};
  }
  \draw[wire] (1.6,0) -- (2.6,0) node[right] {$y = a\oplus b$};
  \draw[wire, FpgaBlue] (0.6,-1.7) -- (0.6,-2.6) node[below, text=FpgaDark] {$a$};
  \draw[wire, FpgaBlue] (1.1,-1.45) -- (1.1,-2.6) node[below, text=FpgaDark] {$b$};
  \node[lbl, align=center, text=FpgaGreen] at (-1.1,2.3) {truth\\table};
  \node[font=\sffamily\small, align=left, anchor=west] at (5.0,0.8) {Substitute \code{0001} --- you get AND,};
  \node[font=\sffamily\small, align=left, anchor=west] at (5.0,0.2) {\code{0111} --- OR, \code{1110} --- NAND\ldots};
  \node[font=\sffamily\small, align=left, anchor=west] at (5.0,-0.4) {The circuit is the same --- only};
  \node[font=\sffamily\small, align=left, anchor=west] at (5.0,-1.0) {the four bits on the left change.};
\end{tikzpicture}
\caption{A 4:1 multiplexer with constants on its inputs computes any function of
two variables (here, exclusive OR)}
\label{fig:mux_lut}
\end{figure}

\begin{ideabox}[Remember this figure]
The four constants on the left are simply a column of values from the truth
table. By changing them we change the function, while the circuit stays the
same. If these bits are stored in a small memory, you get ``programmable''
logic.

\lookahead{a multiplexer with a memory like this is called a LUT; an FPGA
consists of thousands of such blocks (Chapter~\ref{ch:inside}).}
\end{ideabox}

\begin{diybox}[Any function on one multiplexer]
\textbf{Task.} \stars{1} Implement $A\oplus B$ on a 4:1 multiplexer as in
Fig.~\ref{fig:mux_lut}. \stars{2} Implement $F=A\oplus B\oplus C$ (problem~5 of
this chapter): $A,B$ go to the select inputs, and the data inputs get constants
and $C$ or $\overline{C}$. \stars{3} Implement the ``2 out of 3'' (majority)
function from Chapter~\ref{ch:comb} on a single block.

\textbf{In Logra.} The Designer, Waveforms, and RTL tabs; the reference is the
example \code{mux-four}.

\textbf{Expected result.} The output table matches the function's truth table
for all combinations; by changing only the constants on the data inputs, you get
AND, OR, NAND without rebuilding the circuit.

\textbf{How to check.} Waveforms: go through all combinations and compare with
the table; RTL: make sure the structure is the same for different functions.
\end{diybox}

\section{The Demultiplexer}

A \term{demultiplexer} (DEMUX) does the opposite: it has one data input $x$ and
$2^n$ outputs. The address $s$ determines which of the outputs receives the
input signal; the other outputs are 0.

\begin{figure}[H]
\centering
\begin{minipage}{0.35\textwidth}
\centering
\begin{tikzpicture}
  \draw[thick, fill=orange!15] (0,1.2) -- (1.6,2) -- (1.6,-2) -- (0,-1.2) -- cycle;
  \node[font=\sffamily\small\bfseries] at (0.8,0) {DEMUX};
  \draw[wire] (-0.8,0) node[left] {$x$} -- (0,0);
  \foreach \i in {0,...,3} {
    \draw[wire] (1.6,1.5-\i) -- (2.4,1.5-\i) node[right] {$y_\i$};
    \node[font=\sffamily\footnotesize, anchor=east] at (1.55,1.5-\i) {\i};
  }
  \draw[wire, FpgaBlue] (0.5,-1.45) -- (0.5,-2.6) node[below, text=FpgaDark] {$s_1$};
  \draw[wire, FpgaBlue] (1.0,-1.7) -- (1.0,-2.6) node[below, text=FpgaDark] {$s_0$};
\end{tikzpicture}
\end{minipage}\hfill
\begin{minipage}{0.6\textwidth}
\centering
\begin{tabular}{cc|cccc}
\toprule
$s_1$ & $s_0$ & $y_0$ & $y_1$ & $y_2$ & $y_3$ \\
\midrule
0 & 0 & $x$ & 0 & 0 & 0 \\
0 & 1 & 0 & $x$ & 0 & 0 \\
1 & 0 & 0 & 0 & $x$ & 0 \\
1 & 1 & 0 & 0 & 0 & $x$ \\
\bottomrule
\end{tabular}

\medskip
$y_0 = x\,\overline{s_1}\,\overline{s_0}, \ \ y_1 = x\,\overline{s_1}\,s_0, \ \ \ldots$
\end{minipage}
\caption{The 1:4 demultiplexer}
\end{figure}

Compare the formulas with the decoder with an enable input
(Fig.~\ref{fig:dc24}): they coincide if the enable input $E$ is renamed $x$.
Therefore \textbf{a demultiplexer is a decoder whose enable input receives the
data signal}. In reference books these two chips are often listed under a single
name.

\begin{diybox}[The demultiplexer]
\textbf{Task.} \stars{1} Build a demultiplexer: one button, four lamps, and two
switches that set the lamp number. \stars{2} Build the same thing from a decoder
with an enable input $E$, connecting the button to $E$.

\textbf{In Logra.} Learn tab, the ``Selection and distribution'' chapter, the
lesson ``Demultiplexer''; the example \code{demux-four}; for the second task,
the lesson ``Decoder''.

\textbf{Expected result.} A pressed button lights only one lamp, the one chosen
by the address; the others stay off.

\textbf{How to check.} Waveforms: at address 10 the output $y_2$ repeats the
button, the rest are 0; RTL: the block and the decoder with $E$ give identical
waveforms.
\end{diybox}

\section{Transmission over a Single Line}

Together, a multiplexer and a demultiplexer let you send several signals over a
\emph{single} wire. The address quickly cycles through the values $0, 1, 2, 3, 0,
1\ldots$, and at each moment the line carries the signal of one of the sources.
On the receiving side, a demultiplexer with the same address distributes them to
its outputs. This is \term{time-division multiplexing}: it is how telephone
lines and computer buses work, and in drones it is how a remote-control receiver
sends the signals of all channels over one wire (the PPM and SBUS protocols).

\begin{figure}[H]
\centering
\begin{tikzpicture}
  \draw[thick, fill=orange!15] (0,1.3) -- (1.2,0.7) -- (1.2,-0.7) -- (0,-1.3) -- cycle;
  \node[font=\sffamily\footnotesize\bfseries] at (0.6,0) {MUX};
  \foreach \i/\t in {0/throttle, 1/roll, 2/pitch, 3/yaw}
    \draw[wire] (-0.6,0.9-\i*0.6) node[left, font=\sffamily\small] {\t} -- (0,0.9-\i*0.6);
  \draw[very thick, FpgaOrange] (1.2,0) -- node[above, font=\sffamily\small, text=FpgaDark] {single line} (5.0,0);
  \begin{scope}[xshift=5cm]
  \draw[thick, fill=orange!15] (0,0.7) -- (1.2,1.3) -- (1.2,-1.3) -- (0,-0.7) -- cycle;
  \node[font=\sffamily\footnotesize\bfseries] at (0.6,0) {DEMUX};
  \foreach \i/\t in {0/throttle, 1/roll, 2/pitch, 3/yaw}
    \draw[wire] (1.2,0.9-\i*0.6) -- (1.8,0.9-\i*0.6) node[right, font=\sffamily\small] {\t};
  \end{scope}
  \node[block, minimum width=20mm, minimum height=7mm] (ctr) at (3.1,-2.1) {Address\\\footnotesize counter};
  \draw[arr, FpgaBlue] (ctr.west) -| node[pos=0.25, below, lbl] {$s$} (0.6,-1.05);
  \draw[arr, FpgaBlue] (ctr.east) -| node[pos=0.25, below, lbl] {$s$} (5.6,-1.05);
\end{tikzpicture}
\caption{Four remote-control channels over a single wire}
\end{figure}

\begin{diybox}[Security Panel: four zones through one multiplexer]
\textbf{Task.} \stars{1} Build a ``4:1 multiplexer for four sensors'':
$x_0\ldots x_3$ are the zone sensors, the address $s_1s_0$ is set by two
switches, and the output shows the state of the selected zone. \stars{2} Add a
demultiplexer and show the number of the selected zone with a separate lamp.
\stars{3} Replace the address switches with a counter so that the zones are
polled in turn (counters are in Chapter~\ref{ch:regs}; for now you can change
the address by hand).

\textbf{In Logra.} The Designer and Waveforms tabs; references are the lesson
``4-to-1 multiplexer'' and the example \code{mux-four}.

\textbf{Expected result.} The multiplexer output equals the sensor of the
selected zone and depends on nothing else.

\textbf{Security Panel project milestone: milestone~4b.} Four zones through one
multiplexer. Acceptance criterion: for all $4\times 2=8$ combinations of ``zone
address --- the sensor of that zone open/closed'' the output equals the sensor
of the selected zone, and changing the sensors of the other zones does not
affect the output (check on the Waveforms tab).
\end{diybox}

\section{Multibit Multiplexers}

Often we need to choose not between bits but between \emph{numbers}. For this,
several multiplexers are placed in parallel with a common address. For example,
an 8-bit 2:1 multiplexer is eight 2:1 multiplexers whose select inputs are
connected together. In schematics such a unit is drawn as a single symbol with
buses.

\begin{figure}[H]
\centering
\begin{tikzpicture}
  \draw[thick, fill=orange!15] (0,1.1) -- (1.1,0.6) -- (1.1,-0.6) -- (0,-1.1) -- cycle;
  \node[font=\sffamily\small] at (0.55,0) {MUX};
  \draw[bus] (-1.4,0.6) node[left, text=FpgaDark] {$A$} -- (0,0.6);
  \draw[bus] (-1.4,-0.6) node[left, text=FpgaDark] {$B$} -- (0,-0.6);
  \draw[bus] (1.1,0) -- (2.4,0) node[right, text=FpgaDark] {$Y$};
  \foreach \x/\y in {-0.8/0.6, -0.8/-0.6, 1.7/0} {
    \draw[thick] (\x-0.1,\y-0.15) -- (\x+0.1,\y+0.15);
    \node[font=\sffamily\footnotesize, above] at (\x,\y+0.1) {8};
  }
  \draw[wire, FpgaBlue] (0.55,-0.85) -- (0.55,-1.5) node[below, text=FpgaDark] {$s$};
  \node[font=\sffamily\small, anchor=west] at (4.0,0.3) {$s=0$: \ $Y = A$};
  \node[font=\sffamily\small, anchor=west] at (4.0,-0.3) {$s=1$: \ $Y = B$};
\end{tikzpicture}
\caption{A multiplexer for 8-bit numbers. The slash with a number shows how many
wires are in the bus}
\end{figure}

\begin{ideabox}
A \textbf{multiplexer} chooses one of $2^n$ inputs by an address. It consists of
an address decoder and AND--OR gates, and it can implement any function if
constants are applied to its inputs. A \textbf{demultiplexer} passes one input
to the chosen output---it is a decoder with a data signal on its enable input.
\end{ideabox}

\begin{diybox}[Choosing between numbers]
\textbf{Task.} \stars{1} Choose one of three 4-bit numbers $a,b,c$ by a 2-bit
code $s$ (when $s=3$, let $c$ be chosen). \stars{2} Describe the choice between
two 8-bit buses with a single line in Code:
\code{assign y = s ? b : a;}, where \code{a} and \code{b} are 8 bits each.
\stars{3} Add a third source and compare with the ``4:1 multiplexer'' block
whose ``Data width'' parameter is~8.

\textbf{In Logra.} Learn tab, the ``Selection and distribution'' chapter, the
lesson ``Choosing from three sources''; the Code tab for the second and third
tasks; the example \code{mux}.

\textbf{Expected result.} One block (one line of code) switches all the bits at
once; the select input $s$ stays one bit wide (or two bits wide).

\textbf{How to check.} Waveforms: apply two different numbers, for example
\code{0xA5} and \code{0x3C}, and toggle $s$---the output should repeat the
chosen number in full; RTL: one multiplexer symbol with buses, as in the figure
above.
\end{diybox}

\section{Verilog Examples}
\label{sec:mux-verilog}

In Verilog a multiplexer is one line or a short \code{case}. The five examples
below are complete modules from the folder \code{code/b5\_mux/}; each has a
testbench \code{tb\_*.v} that compares the versions with each other and with the
expected value and prints \code{OK}. Running them in Logra: the \emph{Code} tab
$\to$ \emph{Synthesize} (\code{Ctrl+Enter}) $\to$ \emph{RTL} $\to$
\emph{Waveforms} (Section~\ref{sec:signals-verilog}).

\subsection{Example 1. A 2:1 Multiplexer in Three Ways}

One function $y=\overline{s}\,a+s\,b$ is written with the operator \code{?:},
with \code{case}, and with the gates of Fig.~\ref{fig:mux21}.

\vfile{\codepath/b5_mux/mux2.v}{code/b5\_mux/mux2.v}

\textbf{How to read it.}
\begin{itemize}
  \item \code{assign y = s ? b : a;} reads ``if $s$ then $b$, else $a$''. The shortest way.
  \item \code{always @(*)} + \code{case (s)} is the same through a selection statement. The output is declared \code{reg}, but the circuit stays combinational. \code{default} guarantees that \code{y} gets a value for every $s$.
  \item \code{mux2\_gates} is the circuit of Fig.~\ref{fig:mux21}: \code{wire ns = \textasciitilde s;} and two ANDs joined by an OR.
\end{itemize}
\textbf{Checking.} \code{tb\_mux2.v} connects all three modules to the same
inputs and sweeps the eight combinations of $s,a,b$: the outputs must agree.
In Logra synthesize each module and compare the \emph{RTL}: the synthesizer
may produce similar circuits from the three forms.

\textbf{Finish it yourself.} Add an enable input \code{en}: when \code{en} $=0$
the output is zero.

\textbf{Sample answer.}
\begin{vcode}
module mux2_en (
    input  a, b, s, en,
    output y
);
    assign y = en ? (s ? b : a) : 1'b0;   // nested ?:
endmodule
\end{vcode}

\subsection{Example 2. A 4:1 Multiplexer: case, Index and Tree}

There are three convenient forms for four inputs. The tree repeats
Fig.~\ref{fig:muxtree}: two multiplexers of the first level and one of the
second.

\vfile{\codepath/b5_mux/mux4.v}{code/b5\_mux/mux4.v}

\textbf{How to read it.}
\begin{itemize}
  \item \code{input [3:0] x} is the four data inputs as one bus, \code{input [1:0] s} is the address.
  \item \code{x[s]} selects a bus bit by number: a multiplexer by itself.
  \item \code{mux4\_tree}: \code{m0} chooses within the pair ($x_0,x_1$) by the \emph{low} address bit, \code{m1} within ($x_2,x_3$), and \code{y} chooses between the pairs by the high bit. The order of the address bits matters: swap \code{s[0]} and \code{s[1]} and inputs 1 and 2 trade places.
\end{itemize}
\textbf{Checking.} \code{tb\_mux4.v} sweeps all 16 values of \code{x} and all 4
addresses and compares the three modules with \code{x[s]}. In Logra: set
\code{x} $=$ \code{0101} and change \code{s}: the output must alternate
$1,0,1,0$.

\textbf{Finish it yourself.} Build an 8:1 multiplexer (three address bits).

\textbf{Sample answer.}
\begin{vcode}
module mux8 (
    input  [7:0] x,
    input  [2:0] s,
    output       y
);
    assign y = x[s];     // an index by a three-bit address: 8 inputs
endmodule
\end{vcode}
Also try building it as a tree of two 4:1 multiplexers and one 2:1.

\subsection{Example 3. A 1:4 Demultiplexer}

The data input \code{x} must reach the output with number \code{s}, and the other
outputs must be zero. The first form uses \code{case}, the second a shift.

\vfile{\codepath/b5_mux/demux4.v}{code/b5\_mux/demux4.v}

\textbf{How to read it.}
\begin{itemize}
  \item \code{y = 4'b0000;} at the start of the block is a \emph{default value}: the chosen bit is overwritten afterwards. This way \code{y} is assigned on every path and there is no latch, even if \code{case} does not list all the values.
  \item \code{y[2] = x;} assigns to a single bit of the bus.
  \item \code{\{3'b000, x\} << s} shifts the number \code{000x} left by $s$ places: if $x=1$, the one lands in bit $s$. This is a decoder with the enable input $x$.
\end{itemize}
\textbf{Checking.} \code{tb\_demux4.v} sweeps $x$ and all addresses and compares
both modules. In Logra: at \code{s} $=10$ the output \code{y[2]} repeats
\code{x}, and the others are zero.

\textbf{Finish it yourself.} Add an enable input \code{en}: when \code{en} $=0$
all outputs are zero.

\textbf{Sample answer.}
\begin{vcode}
module demux4_en (
    input        x, en,
    input  [1:0] s,
    output [3:0] y
);
    assign y = en ? ({3'b000, x} << s) : 4'b0000;
endmodule
\end{vcode}

\subsection{Example 4. A Multiplexer as a Universal Gate}

XOR comes out of a 2:1 multiplexer if $b$ and $\overline{b}$ are applied to the
data inputs and $a$ is used as the address. Any function of two variables is one
4:1 multiplexer with four constants; in Verilog the constants are collected in a
parameter.

\vfile{\codepath/b5_mux/mux_logic.v}{code/b5\_mux/mux\_logic.v}

\textbf{How to read it.}
\begin{itemize}
  \item \code{a ? \textasciitilde b : b}: at $a=0$ the output is $b$, at $a=1$ it is $\overline b$: exclusive OR.
  \item \code{parameter [3:0] TT} is a truth table column as a module parameter; \code{TT[\{a,b\}]} selects a bit by the two-bit address. The same module with \code{TT = 4'b1000} gives AND, with \code{4'b1110} OR, with \code{4'b0110} exclusive OR. This is the LUT of Fig.~\ref{fig:mux_lut}.
  \item \code{mux\_lut2 \#(.TT(4'b1000)) u (\ldots)} is an instantiation of the module with a given parameter.
\end{itemize}
\textbf{Checking.} \code{tb\_mux\_logic.v} checks XOR, AND, OR, and then creates
16 instances at once with \code{TT} from 0 to 15 and makes sure that for every
row of the table the output equals the corresponding bit of the parameter. In
Logra synthesize \code{xor\_on\_mux} and compare the \emph{RTL} with
Fig.~\ref{fig:mux_lut}.

\textbf{Finish it yourself.} Implement the ``2 out of 3'' function
$F=ab+ac+bc$ on a single 4:1 multiplexer (the address is $\{a,b\}$, the data
inputs are constants and $c$).

\textbf{Sample answer.}
\begin{vcode}
module maj3_mux (
    input  a, b, c,
    output y
);
    // d[0] at ab=00 -> 0;  d[1] at 01 -> c;  d[2] at 10 -> c;  d[3] at 11 -> 1
    wire [3:0] d = {1'b1, c, c, 1'b0};
    assign y = d[{a, b}];
endmodule
\end{vcode}

\subsection{Example 5. Multiplexers on 8-Bit Buses}

Choosing between numbers is written the same way as choosing between bits: the
width of the buses changes nothing.

\vfile{\codepath/b5_mux/mux_bus8.v}{code/b5\_mux/mux\_bus8.v}

\textbf{How to read it.} \code{y = s ? b : a;} with 8-bit \code{a} and \code{b}
switches all eight bits at once: the synthesizer places eight multiplexers with
a common address. The 4:1 version is the same \code{case}, but with a two-bit
address and buses instead of bits; \code{default} covers $s=3$.

\textbf{Checking.} \code{tb\_mux\_bus8.v} applies the numbers \code{0xA5},
\code{0x3C}, \code{0xFF}, \code{0x00} and checks that the output repeats the
chosen number entirely. In Logra set \code{a} $=$ \code{A5}, \code{b} $=$
\code{3C} and toggle \code{s}; on \emph{RTL} this is one multiplexer symbol with
buses, as in the figure ``A multiplexer for 8-bit numbers''.

\textbf{Finish it yourself.} Choose one of three 4-bit numbers $a,b,c$ by the
two-bit code $s$ (at $s=3$ choose $c$), the task from ``Do it yourself''.

\textbf{Sample answer.}
\begin{vcode}
module mux3_bus4 (
    input  [3:0] a, b, c,
    input  [1:0] s,
    output reg [3:0] y
);
    always @(*) begin
        case (s)
            2'd0:    y = a;
            2'd1:    y = b;
            default: y = c;    // s = 2 and s = 3
        endcase
    end
endmodule
\end{vcode}

\section{Common Mistakes}

\begin{warnbox}[Common mistakes]
\begin{itemize}
  \item \textbf{Mixing up the order of the address bits.} In the code $s_1s_0=01$ the high bit is $s_1$, so $x_1$ is selected, not $x_2$. If you swap the wires $s_1$ and $s_0$ in the Designer, zones 1 and 2 trade places, while the circuit still looks like it ``works''.
  \item \textbf{Thinking there are $2^n$ select inputs.} A $2^n{:}1$ multiplexer has $n$ select inputs: a $4{:}1$ has two, not four, and an $8{:}1$ has three.
  \item \textbf{Expecting a demultiplexer to ``remember'' a value.} It has no memory: unselected outputs are always 0, and when the address changes, the previous output goes dark immediately.
  \item \textbf{Letting the addresses at the transmitter and receiver drift apart.} In single-line transmission the multiplexer and the demultiplexer must receive \emph{the same} address at the same moment; even a one-step shift sends the channels to the wrong outputs.
  \item \textbf{Verilog: \code{case} without \code{default} (a latch).}
    \code{always @(*) case (s) 2'd0: y = a; 2'd1: y = b; endcase} leaves $y$
    undefined at $s=2,3$, and the synthesizer inserts a latch that ``remembers''
    the previous value. Either add \code{default} or set a default value at the
    start of the block (\code{y = 1'b0;}).
  \item \textbf{Verilog: an incomplete sensitivity list.}
    \code{always @(a or b)} without \code{s} does not react to an address change
    in simulation, although the synthesized circuit does. Write
    \code{always @(*)}.
  \item \textbf{Verilog: \code{<=} in a combinational block.} In
    \code{always @(*)} use the blocking \code{=}, in the clocked
    \code{always @(posedge clk)} the non-blocking \code{<=}. Mixing them gives a
    mismatch between simulation and the board.
\end{itemize}
\end{warnbox}

\section{Questions and Problems}
\label{sec:tasks-mux}

Answers are collected in Appendix~\ref{app:answers}.

\begin{taskbox}[Chapter problems]
\begin{enumerate}
  \item \stars{1}~How many select inputs does an $8:1$ multiplexer have? And a $16:1$?
  \item \stars{1}~How many $2:1$ multiplexers are needed to build an $8:1$ multiplexer?
  \item \stars{1}~How many $2:1$ multiplexers are needed to choose between two 8-bit buses?
  \item \stars{1}~A $1:4$ demultiplexer has data equal to 1 at address 10. What are the values of the four outputs?
  \item \stars{2}~Implement the function $F = A\oplus B\oplus C$ on a single $4:1$ multiplexer, applying $A$ and $B$ to the select inputs. What must be applied to the data inputs?
  \item \stars{2}~Why is the multiplexer called a universal gate? Which four constants on the inputs $D_0\ldots D_3$ (at address $AB$) give the AND and OR functions?
  \item \stars{2}~\emph{Find the error.} A student implements $A\oplus B$ on a $4:1$ multiplexer with address $AB$ and applies the constants $0,1,1,1$ to $D_0,D_1,D_2,D_3$. Which function did they actually get, and which constants are needed?
  \item \stars{2}~\emph{Find the error.} In a single-line transmission system (the section ``Transmission over a Single Line''), the receiver's address counter is started one clock cycle later than the transmitter's counter. What will appear on the outputs ``throttle'', ``roll'', ``pitch'', ``yaw''?
  \item \stars{3}~\textbf{Project: Security Panel milestone 4b.} Four zone sensors through one $4:1$ multiplexer, the zone address set by two switches; a demultiplexer shows the zone number with a lamp. \emph{Acceptance criterion:} at each of the four addresses the output repeats the sensor of exactly that zone, and changing the other zones has no effect; only one number lamp is lit (Waveforms tab).
  \item \stars{3}~\textbf{Project: ``Three sources''.} Choose one of three 4-bit numbers $a,b,c$ by a 2-bit code $s$ (when $s=3$ --- $c$). \emph{Acceptance criterion:} for $s=0,1,2,3$ the output equals $a,b,c,c$ (Waveforms tab, any numbers in the example).
  \item \stars{3}~\textbf{Project: ``Majority vote''.} Implement the ``2 out of 3'' function $F=AB+AC+BC$ on a single $4:1$ multiplexer. \emph{Acceptance criterion:} the output table matches the truth table from Chapter~\ref{ch:comb} for all 8 combinations.
\end{enumerate}
\end{taskbox}

\section{What's Next}

We have learned to select and distribute: to get one signal out of many and
back again. But we can select not only sensors but also the \emph{results of
computations}: let the sum, the difference, and the logical AND be computed
first, and let a multiplexer pass the needed answer according to an operation
code. A question remains: how does a circuit made of gates \emph{calculate} at
all? In the next chapter (Chapter~\ref{ch:arith}) we will build an adder, a
subtractor, and a comparator, and from them and the multiplexer of this chapter
we will assemble an arithmetic logic unit.
